\section{Gamma reciprocity, polygamma sums, and a normalized primitive}
\label{asj:sec:gamma}

\subsection{A convergent nonlinear part of the first Laurent jet}
For the multiplicative Hurwitz lattice define
\begin{equation}\label{asj:eq:Edef}
 E_{\boldsymbol b}(a)=\sum_{\boldsymbol n}
        \left\{\frac a{\lambda_{\boldsymbol n}}
                    -\log\left(1+\frac a{\lambda_{\boldsymbol n}}\right)\right\},
 \qquad a> -\lambda_*.
\end{equation}
Every brace must be kept grouped. Its tail is $O(\lambda_{\boldsymbol n}^{-2})$
on compact shift sets, and $\sum\lambda_{\boldsymbol n}^{-2}
=\prod_j\zeta(2,b_j)<\infty$. Thus this is an ordinary, locally normally
convergent series, not a finite-part sum.

\begin{proposition}[First-jet decomposition]\label{asj:thm:firstjet}
For $F_a$ of Section~\ref{asj:sec:hurwitz},
\begin{equation}\label{asj:eq:firstjet}
 \boxed{\displaystyle \calJ_1F_a=\D'(0)-a d_0+E_{\boldsymbol b}(a).}
\end{equation}
\end{proposition}
\begin{proof}
For $|a|<\lambda_*$, apply \eqref{asj:eq:alljets} with $P=1$, one exponent,
and $r=1$. The resonant term is $-a d_0$. For $\ell\ge2$,
$[s](s)_\ell/\ell!=1/\ell$, so the remaining sum is
\[
 \sum_{\ell\ge2}\frac{(-a)^\ell}{\ell}\D(\ell)
       =\sum_{\boldsymbol n}\left\{a/\lambda_{\boldsymbol n}
                  -\log(1+a/\lambda_{\boldsymbol n})\right\}.
\]
Absolute convergence justifies the interchange. Both sides continue
locally holomorphically throughout the slit shift plane
$\C\setminus(-\infty,-\lambda_*]$, using finite heads on the spectral
side and the grouped series on the other side. The identity theorem,
or continuation along the real interval, gives every stated $a$.
\end{proof}

The entire principal part at the spectral origin is affine in $a$ by
\eqref{asj:eq:Fpolar}. All nonlinearity of the first Laurent coefficient is
therefore the convergent function $E_{\boldsymbol b}$.

\subsection{Two independent Gamma representations of the same function}
For two parameters write $E_{b,c}=E_{(b,c)}$.

\begin{theorem}[Centered Gamma reciprocity]\label{asj:thm:gamma}
For $b,c>0$ and $a>-bc$,
\begin{align}\label{asj:eq:gamma-reciprocity}
 E_{b,c}(a)
 &=\sum_{n=0}^{\infty}\left\{
       \log\Gamma\left(c+\frac a{n+b}\right)-\log\Gamma(c)
                      -\frac{a\psi(c)}{n+b}\right\}\notag\\
 &=\sum_{m=0}^{\infty}\left\{
       \log\Gamma\left(b+\frac a{m+c}\right)-\log\Gamma(b)
                      -\frac{a\psi(b)}{m+c}\right\}.
\end{align}
Both series converge absolutely and locally uniformly in $a$. On
$|a|<bc$ they also equal
\begin{equation}\label{asj:eq:gamma-zeta}
 E_{b,c}(a)=\sum_{\ell\ge2}
                 \frac{(-a)^\ell}{\ell}\zeta(\ell,b)\zeta(\ell,c).
\end{equation}
\end{theorem}
\begin{proof}
The classical centered Gamma product gives, for $c>0$ and $t>-c$,
\begin{equation}\label{asj:eq:centeredgamma}
 \log\Gamma(c+t)-\log\Gamma(c)-t\psi(c)
    =\sum_{m\ge0}\left\{\frac{t}{m+c}
                         -\log\left(1+\frac t{m+c}\right)\right\}.
\end{equation}
One may derive it without choosing a product normalization: both sides
vanish with their first derivative at $t=0$, and their second derivatives
are the same convergent trigamma series. Set $t=a/(n+b)$ and sum over
$n$. All double summands are absolutely summable, since their tails are
bounded by a constant times $(n+b)^{-2}(m+c)^{-2}$ and their finitely many
small factors stay away from zero on compact allowed shift sets.
Fubini's theorem gives $E_{b,c}$. Interchanging $m$ and $n$ proves the
second representation. Expanding the logarithm under $|a|<bc$ proves
\eqref{asj:eq:gamma-zeta}.
\end{proof}

For $b=c=1$, the first Laurent coefficient becomes the explicit
Gamma--Stieltjes identity
\begin{equation}\label{asj:eq:divisor-gamma}
 \calJ_1F_a=\frac12\log(2\pi)-a(\gamma^2-2\gamma_1)
      +\sum_{n\ge1}\left\{\log\Gamma(1+a/n)+\frac{\gamma a}{n}\right\}.
\end{equation}
Here $\D'(0)=2\zeta(0)\zeta'(0)=\tfrac12\log(2\pi)$ and
$d_0=\gamma^2-2\gamma_1$. At the same time,
\[
 \Res_{s=0}F_a=-a,\qquad [s^0]F_a=\frac14-2\gamma a.
\]
The Gamma series in \eqref{asj:eq:divisor-gamma} cannot make that nonzero pole
 disappear. It evaluates the correctly specified Laurent coefficient.

\subsection{Reciprocal digamma and all polygamma derivatives}
\begin{corollary}[All argument derivatives]\label{asj:thm:polygamma}
For $b,c>0$, $a>-bc$,
\begin{align}\label{asj:eq:digamma-reciprocity}
 E_{b,c}'(a)
 &=\sum_{n\ge0}\frac{\psi(c+a/(n+b))-\psi(c)}{n+b}\notag\\
 &=\sum_{m\ge0}\frac{\psi(b+a/(m+c))-\psi(b)}{m+c}.
\end{align}
For every integer $r\ge2$,
\begin{align}\label{asj:eq:polygamma-reciprocity}
 E_{b,c}^{(r)}(a)
 &=\sum_{n\ge0}\frac{\psi^{(r-1)}(c+a/(n+b))}{(n+b)^r}\notag\\
 &=\sum_{m\ge0}\frac{\psi^{(r-1)}(b+a/(m+c))}{(m+c)^r}\notag\\
 &=(-1)^r(r-1)!\sum_{n,m\ge0}
                       \bigl((n+b)(m+c)+a\bigr)^{-r}.
\end{align}
At $a=0$ this is $(-1)^r(r-1)!\zeta(r,b)\zeta(r,c)$.
\end{corollary}
\begin{proof}
Differentiate the grouped normally convergent series on compact subsets
of the allowed shift domain. The $r$th derivative for $r\ge2$ of one
grouped summand in \eqref{asj:eq:Edef} is
$(-1)^r(r-1)!(\lambda+a)^{-r}$. The Gamma representation gives the
first two forms. The conventional polygamma series, or differentiation
of \eqref{asj:eq:centeredgamma}, gives the last one. At $a=0$ absolute
convergence factors the double series.
\end{proof}

For more than two Hurwitz factors, summing over one factor gives an
analogous centered Gamma series over the remaining $(h-1)$ indices.
The choice of the factor summed first is arbitrary, producing $h$
compatible Gamma representations. This is a constructive representation,
not a reduction claim into finitely many Gamma values.

\subsection{An anchored primitive in Hurwitz-zeta derivatives}
Define the classical primitive
\begin{equation}\label{asj:eq:Kdef}
 K(z)=\zeta'(-1,z)+\frac{z-z^2}{2}+\frac z2\log(2\pi),\qquad z>0.
\end{equation}
The identities $\partial_z\zeta(s,z)=-s\zeta(s+1,z)$ and
$\zeta'(0,z)=\log\Gamma(z)-\tfrac12\log(2\pi)$ imply
$K'(z)=\log\Gamma(z)$ \cite{dlmf-hurwitz}. This agrees with the
primitive already used in the manuscript \cite{repo}.

\begin{theorem}[Reciprocal normalized primitive]\label{asj:thm:primitive}
Let $I_{b,c}(a)=\int_0^a E_{b,c}(t)\dd t$, with oriented integration
when $a<0$. Then
\begin{align}\label{asj:eq:primitive-gamma}
 I_{b,c}(a)=\sum_{n\ge0}\Bigg\{&(n+b)
              \left[K\left(c+\frac a{n+b}\right)-K(c)\right]\notag\\
              &-a\log\Gamma(c)-\frac{a^2\psi(c)}{2(n+b)}\Bigg\}.
\end{align}
The same expression with $b$ and $c$ exchanged is equal to it. Equivalently,
\begin{align}\label{asj:eq:primitive-lattice}
 I_{b,c}(a)=\sum_{n,m\ge0}\Bigg\{&\frac{a^2}{2\lambda_{n,m}}
        -(\lambda_{n,m}+a)
                 \log\left(1+\frac a{\lambda_{n,m}}\right)+a\Bigg\},
\end{align}
where $\lambda_{n,m}=(n+b)(m+c)$. For $|a|<bc$,
\begin{equation}\label{asj:eq:primitive-power}
 I_{b,c}(a)=\sum_{\ell\ge2}\frac{(-1)^\ell a^{\ell+1}}{\ell(\ell+1)}
                                   \zeta(\ell,b)\zeta(\ell,c).
\end{equation}
\end{theorem}
\begin{proof}
Integrate the grouped summands of Theorem~\ref{asj:thm:gamma} along the
compact segment from $0$ to $a$. Uniform convergence permits the
interchange. Substitution $u=c+t/(n+b)$ gives the $K$ term, and the
two subtractions integrate to the other two terms in
\eqref{asj:eq:primitive-gamma}. Integration of \eqref{asj:eq:Edef} gives
\eqref{asj:eq:primitive-lattice}. Its grouped tail is $O(\lambda_{n,m}^{-2})$.
Both forms vanish at $a=0$, fixing the integration constant. Integrating
\eqref{asj:eq:gamma-zeta} proves \eqref{asj:eq:primitive-power}.
\end{proof}

In \eqref{asj:eq:primitive-gamma}, the three large pieces must not be summed
separately. The cancellations are part of the identity's definition.
This is the same normalization issue that occurs in negative polygamma,
but here both the base point and every convergent subtraction are explicit.
